%!TEX root = main.tex
% ==============================================================================
% İSTANBUL ÜNİVERSİTESİ SOSYAL BİLİMLER ENSTİTÜSÜ
% BÖLÜM 3 (Madde 16/c, d ve Madde 17/b)
% Madde 16/d: Giriş kısmı hariç her tez en az üç bölümden oluşmalıdır.
% ==============================================================================
\chapter{TARTIŞMA VE ELEŞTİREL DEĞERLENDİRME}

\section{Bulguların Alanyazın Işığında Değerlendirilmesi}

Elde edilen bulgular, daha önce yapılmış kuramsal ve ampirik çalışmalarla karşılaştırmalı olarak tartışılmaktadır.\footnote{Her bölüm başında olduğu gibi dipnot sayacı 1'den başlar.} Hukuk felsefesi ve sosyolojisinde iktidar ilişkilerinin analizi için Hirsch'in değerlendirmeleri dikkate değerdir \parencite{hirsch_iktidar}.

\subsection{Alt Boyutlar ve Kuramsal Yansımalar}

Bu alt başlıkta, kuram ile uygulama arasındaki uyum veya farklılaşma noktaları irdelenmektedir.

\subsubsection{Birinci Alt Boyut: Yapısal Faktörler}

Yapısal unsurlar, kurumsal mekanizmaların işleyişi üzerinde belirleyici role sahiptir.

\subsubsection{İkinci Alt Boyut: İşlevsel Dinamikler}

İşlevsel dinamikler ise kurumsal dönüşümlerin hızını ve yönünü tayin etmektedir.

\section{Çalışmanın Sınırları ve Yöntemsel Kısıtlar}

Araştırma sürecinde ulaşılan kaynaklar ve veri toplama imkânları çerçevesinde ortaya çıkan yöntemsel kısıtlılıklar bu kısımda açıklanmıştır.



